% AI4Pain 2026 Challenge -- ACII 2026 PAAIn Workshop submission
% Phase 7 scaffold: structural skeleton only. The PAAIn workshop CFP and
% LaTeX template have not been released yet. Once the organizers publish the
% official template, swap the documentclass and styling to match. The
% sectional skeleton below should remain valid regardless.

\documentclass[conference]{IEEEtran}
\IEEEoverridecommandlockouts

% =============================================================================
% Packages
% =============================================================================
\usepackage{cite}
\usepackage{amsmath,amssymb,amsfonts}
\usepackage{algorithmic}
\usepackage{graphicx}
\usepackage{textcomp}
\usepackage{xcolor}
\usepackage{booktabs}
\usepackage{hyperref}

% =============================================================================
% Metadata (placeholders, fill at submission time)
% =============================================================================
\title{Multi-Signal Pain Localization for the AI4Pain 2026 Challenge}

\author{
\IEEEauthorblockN{Vignan Kamarthi}
\IEEEauthorblockA{
Northeastern University \\
Boston, MA, USA \\
\textit{kamarthi.v@northeastern.edu}
}
}

% =============================================================================
\begin{document}
\maketitle

% =============================================================================
% Abstract
% =============================================================================
\begin{abstract}
[TBD: 150--250 word abstract.]
We present a 3-class pain localization system for the AI4Pain 2026 Challenge,
classifying physiological recordings as No Pain, Hand Pain, or Arm Pain from
four signals (BVP, EDA, RESP, SpO2) collected over 65 participants
(41 / 12 / 12 train / validation / test split). Our approach combines a
classical machine learning track using Catch22, entropy, and statistical
features extracted via a Rust pipeline with three deep learning architectures
operating on raw multi-channel windows: a ResNet-1D baseline, a current-SOTA
multimodal architecture (Phase 4 research), and a novel framework
(Phase 5 research) that exploits [physiological prior to be filled in].
[Headline result placeholder.] Our approach is constrained to operate without
external data and submits five label predictions to the challenge organizers.
\end{abstract}

\begin{IEEEkeywords}
pain assessment, physiological signals, multimodal fusion, biosignal classification, ACII 2026
\end{IEEEkeywords}

% =============================================================================
% 1. Introduction
% =============================================================================
\section{Introduction}
[TBD] Motivation: clinical relevance of automated pain localization,
contrast with binary or intensity-based pain classification, why
multimodal physiology beats single-channel approaches.

Our contributions:
\begin{itemize}
  \item [Contribution 1: e.g., reproducible Rust feature extractor adapted for the four AI4Pain signals]
  \item [Contribution 2: e.g., principled comparison of classical ML and deep learning under matching subject splits and ablations]
  \item [Contribution 3: novel architecture (Phase 5) and what it captures about pain localization that prior SOTA does not]
\end{itemize}

% =============================================================================
% 2. Related Work
% =============================================================================
\section{Related Work}
[Populated by Phase 4 literature review (research/NN2\_SOTA\_RESEARCH.md).]
\subsection{Pain Classification from Physiological Signals}
\subsection{Multimodal Fusion for Biosignals}
\subsection{Pain Localization vs.\ Intensity}

% =============================================================================
% 3. Dataset
% =============================================================================
\section{Dataset and Preprocessing}
The AI4Pain 2026 Challenge provides recordings from 65 participants
[citation TBD] with a fixed subject-level split of 41 train / 12 validation /
12 test subjects. Each trial contains four synchronized physiological
channels:
\begin{itemize}
  \item BVP (blood volume pulse)
  \item EDA (electrodermal activity)
  \item RESP (respiration)
  \item SpO2 (blood oxygen saturation)
\end{itemize}
Trial labels are one of \texttt{No Pain (NP)}, \texttt{Hand Pain (HP)}, or
\texttt{Arm Pain (AP)}.

[TBD: per-signal sampling rates, recording duration, label distribution table.
These fields are confirmed at Runtime HIP RT-A when the labeled release lands.]

\subsection{Subject-level Splits}
We strictly follow the challenge's fixed split. No subject appears in
multiple splits and no preprocessing statistic is fit on subjects outside the
training split.

\subsection{Ablation Configurations}
We evaluate three signal subsets to test the incremental value of each modality:

\begin{table}[h]
\centering
\caption{Ablation configurations}
\begin{tabular}{lcc}
\toprule
Config & Signals & Total Features (classical) \\
\midrule
\texttt{bvp\_eda} & BVP + EDA & 80 \\
\texttt{bvp\_eda\_resp} & BVP + EDA + RESP & 120 \\
\texttt{all\_four} & BVP + EDA + RESP + SpO2 & 160 \\
\bottomrule
\end{tabular}
\end{table}

% =============================================================================
% 4. Methods
% =============================================================================
\section{Methods}

\subsection{Classical Machine Learning Track}
We extract 40 features per signal via a Rust feature extractor adapted from
prior work on cuffless blood pressure inference: 22 Catch22 canonical
time-series features, 10 entropy measures (8 ordinal-pattern entropies plus
sample and approximate entropy), and 8 statistical features (mean, median,
std, skewness, kurtosis, RMS, min, max). After per-column standardization fit
on the training subjects only, we tune four classifiers via Optuna with
stratified subject-disjoint 5-fold cross-validation, optimizing balanced
accuracy:
\begin{itemize}
  \item Logistic Regression (linear baseline)
  \item Random Forest Classifier
  \item XGBoost Classifier
  \item LightGBM Classifier
\end{itemize}

\subsection{Deep Learning Track}
We train three architectures on raw multi-channel signal windows
[window length and stride confirmed at RT-E]:

\subsubsection{NN \#1: ResNet-1D Baseline}
A 1D ResNet adapted from prior cardiovascular signal work, with the
classification head reshaped from regression output (1 unit) to 3-class
logits. Per-channel z-scoring uses statistics fit on training subjects only.

\subsubsection{NN \#2: Current SOTA}
[Filled in by Phase 4 research; cite the selected architecture and rationale.]

\subsubsection{NN \#3: Novel Framework}
[Filled in by Phase 5 research; this is the paper's central contribution.
Block diagram and ablation plan to be inserted as Figure~\ref{fig:nn3}.]

\subsection{Training Protocol}
All deep models use \texttt{nn.CrossEntropyLoss} with class weights derived
from the training label distribution, AdamW optimization with weight decay,
\texttt{ReduceLROnPlateau} on validation balanced accuracy, gradient clipping
at norm 1.0, per-epoch checkpointing, and early stopping (patience 20). Best
checkpoint is selected by validation balanced accuracy.

% =============================================================================
% 5. Experiments
% =============================================================================
\section{Experiments}

\subsection{Evaluation Metrics}
We report balanced accuracy (primary), macro-F1, per-class precision / recall
for No Pain, Hand Pain, and Arm Pain, full $3 \times 3$ confusion matrices,
and one-vs-rest AUC. We never report a single aggregate score in isolation.

\subsection{Hardware}
1$\times$NVIDIA H200 GPU on the Northeastern Explorer cluster, 8-hour wall
time per job, 12 GB RAM. No model training exceeds the wall time budget.

% =============================================================================
% 6. Results
% =============================================================================
\section{Results}

\begin{table}[h]
\centering
\caption{Validation results across ablation configurations [TBD].}
\begin{tabular}{llcccc}
\toprule
Config & Model & Bal. Acc. & Macro-F1 & AUC-OVR & Notes \\
\midrule
\texttt{bvp\_eda} & Logistic & -- & -- & -- & -- \\
\texttt{bvp\_eda} & RF & -- & -- & -- & -- \\
\texttt{bvp\_eda} & XGB & -- & -- & -- & -- \\
\texttt{bvp\_eda} & LGBM & -- & -- & -- & -- \\
\texttt{bvp\_eda} & ResNet-1D & -- & -- & -- & -- \\
\texttt{bvp\_eda} & NN \#2 & -- & -- & -- & -- \\
\texttt{bvp\_eda} & NN \#3 & -- & -- & -- & -- \\
\bottomrule
\end{tabular}
\end{table}

\subsection{Test Set Submissions}
The challenge limits each team to five test-set label submissions. We track
each submission with rationale and post-hoc result in our internal
\texttt{submission\_tracker.md} (see Section TBD).

% =============================================================================
% 7. Discussion
% =============================================================================
\section{Discussion}
[TBD]

% =============================================================================
% 8. Conclusion
% =============================================================================
\section{Conclusion}
[TBD]

% =============================================================================
% Acknowledgments
% =============================================================================
\section*{Acknowledgments}
Compute provided by the Northeastern University Explorer cluster.

% =============================================================================
% References
% =============================================================================
\bibliographystyle{IEEEtran}
\bibliography{bibliography}

\end{document}
